\documentclass[11pt,a4paper]{article}
\usepackage[utf8]{inputenc}
\usepackage[vietnamese]{babel}
\usepackage{amsmath,amssymb}
\usepackage{geometry}
\geometry{a4paper, total={180mm,260mm}, left=15mm, right=15mm, top=18mm, bottom=18mm}
\usepackage{tikz}
\usetikzlibrary{shapes.geometric, arrows.meta, positioning, calc}
\usepackage{tcolorbox}
\tcbuselibrary{skins,breakable}
\usepackage{xcolor}
\usepackage{multicol}
\usepackage{array}
\usepackage{fontawesome}
\usepackage{fancyhdr}

\pagestyle{fancy}
\fancyhf{}
\renewcommand{\headrulewidth}{0.5pt}
\renewcommand{\footrulewidth}{0.5pt}
\fancyhead[L]{\small\bfseries\color{teal}\faBook\ SỔ TAY KIẾN THỨC TOÁN 3}
\fancyhead[R]{\small\bfseries\color{teal}Lớp Toán Thầy Lĩnh -- loptoanthaylinh.vn}
\fancyfoot[C]{\small\color{gray}-- Sổ tay kiến thức: Bài 65. Luyện tập chung: Phép cộng và phép trừ phạm vi 100 000 --}

\definecolor{primary}{HTML}{005F73}
\definecolor{secondary}{HTML}{0A9396}
\definecolor{accent}{HTML}{EE9B00}
\definecolor{danger}{HTML}{AE2012}
\definecolor{boxbg}{HTML}{F4F6F8}

\begin{document}

% --- BANNER TIÊU ĐỀ ---
\begin{tcolorbox}[
  enhanced,
  colback=white,
  colframe=primary,
  arc=3mm,
  boxrule=1.2pt,
  left=4mm, right=4mm, top=2.5mm, bottom=2.5mm,
  drop shadow=black!10!white
]
  \centering
  {\Large\bfseries\color{primary}\faBookmark\ BÀI 65: LUYỆN TẬP CHUNG: CỘNG TRỪ PHẠM VI 100 000}\\[1.5mm]
  {\small\itshape\color{black!75}Đại tổng kết cộng trừ 5 chữ số, tính giá trị biểu thức số, tìm x và giải bài toán kinh tế thực tế trên 1 trang A4}
\end{tcolorbox}

\vspace{-2mm}


\noindent
\begin{minipage}[t]{0.49\textwidth}\vspace{0pt}
\begin{tcolorbox}[
  enhanced,
  colback=blue!4!white,
  colframe=blue!75!black,
  arc=2mm, boxrule=0.8pt,
  left=2mm, right=2mm, top=1.8mm, bottom=1.8mm,
  title={\footnotesize\bfseries\color{white}\faCalculator\ 1. Kĩ thuật tính toán tổng hợp}
]
{\footnotesize
\begin{itemize}\setlength{\itemsep}{1.5pt}\setlength{\leftskip}{-2.5mm}
  \item Đặt tính thẳng cột tuyệt đối giữa các hàng.
  \item Thực hiện tính từ phải sang trái.
  \item Thuộc lòng các cặp số tròn chục nghìn:
  $$10\,000 + 90\,000 = 20\,000 + 80\,000 = \dots = 100\,000$$
  \item Thứ tự biểu thức: có ngoặc làm trong ngoặc trước, chỉ có cộng trừ làm từ trái sang phải.
\end{itemize}
}
\end{tcolorbox}
\end{minipage}
\hfill
\begin{minipage}[t]{0.49\textwidth}\vspace{0pt}
\begin{tcolorbox}[
  enhanced,
  colback=teal!4!white,
  colframe=teal!75!black,
  arc=2mm, boxrule=0.8pt,
  left=2mm, right=2mm, top=1.8mm, bottom=1.8mm,
  title={\footnotesize\bfseries\color{white}\faSliders\ 2. Giải bài toán có lời văn 2 bước}
]
{\footnotesize
\begin{itemize}\setlength{\itemsep}{1.5pt}\setlength{\leftskip}{-2.5mm}
  \item \textbf{Bài toán thu chi / bán hàng:}\\
  Bước 1: Tính số tiền hoặc sản phẩm đã bán/nhập thêm.\\
  Bước 2: Tính số tiền hoặc sản phẩm còn lại.
  \item \textbf{Bài toán quãng đường / sản lượng:}\\
  Cộng tổng hai giai đoạn hoặc tìm phần chênh lệch (nhiều hơn / ít hơn).
\end{itemize}
}
\end{tcolorbox}
\end{minipage}

\vspace{-2.5mm}
% --- KHỐI 3: SƠ ĐỒ TƯ DUY TIKZ ---
\noindent
\begin{tcolorbox}[
  enhanced,
  colback=orange!4!white,
  colframe=orange!80!black,
  arc=2mm, boxrule=0.8pt,
  left=2mm, right=2mm, top=1.5mm, bottom=1.5mm,
  title={\footnotesize\bfseries\color{white}\faSitemap\ SƠ ĐỒ TƯ DUY: BẢN ĐỒ CỘNG TRỪ TRONG PHẠM VI 100 000}
]
\centering
\begin{tikzpicture}[
  box/.style={draw, rounded corners=2pt, align=center, font=\scriptsize\bfseries, inner sep=2.5pt, minimum height=0.75cm},
  arr/.style={-{Stealth[length=2mm]}, thick}
]
  \node[box, fill=purple!15, draw=purple!80!black, text width=2.6cm] (c1) at (0, 0) {ĐẶT TÍNH\\Thẳng cột 5 hàng\\Cộng/Trừ phải qua trái};
  \node[box, fill=teal!15, draw=teal!80!black, text width=2.6cm] (c2) at (3.5, 0) {TÍNH NHẨM\\Số tròn nghìn\\Số tròn chục nghìn};
  \node[box, fill=blue!12, draw=blue!80!black, text width=2.6cm] (c3) at (7.0, 0) {TÌM THÀNH PHẦN\\Tìm số hạng, số bị trừ\\Tìm số trừ trong phép tính};
  \node[box, fill=orange!15, draw=orange!80!black, text width=2.6cm] (c4) at (10.5, 0) {TOÁN THỰC TẾ\\Thu chi, bán hàng\\Khoảng cách địa lí};

  \draw[arr, purple!80!black] (c1) -- (c2);
  \draw[arr, teal!80!black] (c2) -- (c3);
  \draw[arr, blue!80!black] (c3) -- (c4);
\end{tikzpicture}
\end{tcolorbox}

\vspace{-2.5mm}
% --- KHỐI 4: VÍ DỤ & MẸO TRÁNH BẪY ---
\noindent
\begin{minipage}[t]{0.49\textwidth}\vspace{0pt}
\begin{tcolorbox}[
  enhanced,
  colback=green!4!white,
  colframe=green!60!black,
  arc=2mm, boxrule=0.8pt,
  left=2mm, right=2mm, top=1.8mm, bottom=1.8mm,
  title={\footnotesize\bfseries\color{white}\faCheckCircle\ Ví dụ thực hành chuẩn mực}
]
{\footnotesize
\textbf{1. Tính nhanh giá trị biểu thức:}\\
$A = 24\,000 + 36\,000 - 15\,000 = 60\,000 - 15\,000 = \mathbf{45\,000}$.\\
$B = 100\,000 - (40\,000 + 25\,000) = 100\,000 - 65\,000 = \mathbf{35\,000}$.\\
\textbf{2. Bài toán thực tế:}\\
Một nông trại thu hoạch được $25\,400\text{ kg}$ thóc vụ chiêm và $32\,600\text{ kg}$ thóc vụ mùa. Nông trại đã bán đi $40\,000\text{ kg}$. Còn lại:\\
Tổng số thóc: $25\,400 + 32\,600 = 58\,000\text{ (kg)}$.\\
Còn lại: $58\,000 - 40\,000 = \mathbf{18\,000\text{ (kg)}}$.
}
\end{tcolorbox}
\end{minipage}
\hfill
\begin{minipage}[t]{0.49\textwidth}\vspace{0pt}
\begin{tcolorbox}[
  enhanced,
  colback=red!4!white,
  colframe=red!70!black,
  arc=2mm, boxrule=0.8pt,
  left=2mm, right=2mm, top=1.8mm, bottom=1.8mm,
  title={\footnotesize\bfseries\color{white}\faExclamationTriangle\ Mẹo tránh bẫy trắc nghiệm}
]
{\footnotesize
\begin{itemize}\setlength{\itemsep}{1.5pt}\setlength{\leftskip}{-2.5mm}
  \item \textbf{Bẫy không tính trong ngoặc trước:}\\
  Ở biểu thức $100\,000 - (40\,000 + 25\,000)$, nhiều bạn bỏ ngoặc và tính $100\,000 - 40\,000 = 60\,000 + 25\,000 = 85\,000$ là \textbf{SAI}!
  \item Bắt buộc thực hiện phép tính trong ngoặc tròn trước.
\end{itemize}
}
\end{tcolorbox}
\end{minipage}


\end{document}
